\subsection*{Terminal gauges and the A.11 axial stress source}
Source (A.11) preserves the pressure increment, but the separate
equation (A.19) fixes $S(\infty,\eta)=0$.
The pressure terminal condition requires the source's prescribed
axis datum (A.21)--(A.23) and the exact pressure-neutral A.11
correction. They must hold for the \emph{same} corrected profile.
Write $d=1-\eta^2$, $A=1/2+h$, and define the terminal defects
$\sigma=S(\infty,\eta)$ and $\pi=\Pi(\infty,\eta)$.
On both A.11 supports $U=M=J=0$, so (4.15)--(4.16) give
\begin{equation}\label{eq:A11-terminal-gauge-full}
 N_s=N_s^{\rm tail}
 +\frac{4h\eta\sigma-d\sigma_\eta}{X}
 +4A\eta\pi-d\pi_\eta .
\end{equation}
Here $N_s^{\rm tail}$ substitutes
$S_{\rm tail}=\frac12\int_X^\infty E^2dx$ and
$\Pi_{\rm tail}=-\frac12\int_X^\infty E^2dx/x$.
The identity holds before terminal normalization and implies
\[
 \frac{|N_s-N_s^{\rm tail}|}{E}
 \le\frac{4h|\sigma|+|\sigma_\eta|}{XE}
 +\frac{4A|\pi|+|\pi_\eta|}{E}.
\]
It explicitly demonstrates why the two A.11 equations alone
do not certify $|N_s|/E$.
If $\sigma=\pi=0$ and differentiation under the convergent
integrals is justified, then the four tail quantities combine
into one source-exact integral:
\begin{equation}\label{eq:A11-single-future-gate}
 \frac{N_s}{E(X,\eta)}
 =\frac1{E(X,\eta)}\int_X^\infty E(x,\eta)^2
 \left[\frac{2h\eta}{X}-\frac{2A\eta}{x}
 +(1-\eta^2)\partial_\eta\log E(x,\eta)
 \left(\frac1x-\frac1X\right)\right]dx .
\end{equation}
Define $D_s=S_{\rm tail}/(XE^2)$ and
$P_s=-\Pi_{\rm tail}/E^2$. Then $0\le P_s\le D_s$.
If the full future logarithmic slope $\ell\le0$, also
$P_s\le1/2$. If $|\partial_\eta\log E|\le\varepsilon$
at every future radius, the integral gives
\[
 \frac{|N_s|}{E}\le E\left[
 4|\eta|(hD_s+AP_s)
 +2(1-\eta^2)\varepsilon(D_s-P_s)\right].
\]
Thus the jointly certified sufficient source inputs
$0<\lambda\le1/20$, $0<h\le\lambda/100$,
$E\le\lambda/4$, $D_s\le233/100$,
$\varepsilon\le1/1000$, $\ell\le0$, and
$\sigma=\pi=0$ yield the exact-rational implication
\begin{equation}\label{eq:A11-terminal-strict-slack}
 |N_s|/E<253\lambda/1000<\lambda/2 .
\end{equation}
For imperfect gauges, the same result holds if the
two displayed gauge-error ratios sum to less than
$247\lambda/1000$. These are conditional inequalities,
not independently source-authenticated enclosures.

To see the unequal conditioning of the two terminal defects,
write $X=X_pe^y$, $f=(1+\eta^2)^{-1}$,
$E=e_b f e^{-(1/2+\lambda)y}G$ with $G>0$,
$\sigma=X_pe_b^2f^2r/\lambda$, and $\pi=e_b^2f^2p$.
The corresponding normalized corrections to $N_s/E$ are
\[
 \frac{e_bf}{\lambda G}e^{-(1/2-\lambda)y}
 \left[4h\eta r-d\left(2(f_\eta/f)r+r_\eta\right)\right],
\quad
 \frac{e_bf}{G}e^{(1/2+\lambda)y}
 \left[4A\eta p-d\left(2(f_\eta/f)p+p_\eta\right)\right].
\]
The $S$-defect coefficient decays with radius;
the pressure-gauge coefficient grows. This establishes
why the exact terminal roots or authenticated derivative
enclosures are required before applying the stress gate.
It does not close the $Q_s$, complete exterior cone,
Appendix B, or proposed unforced Theorem~1.1.